\documentclass[11pt,a4paper]{article}
\usepackage[margin=2.2cm]{geometry}
\usepackage{amsmath}
\usepackage{graphicx}
\usepackage{booktabs}
\usepackage{subcaption}
\usepackage{url}
\usepackage[round]{natbib}
\usepackage[hidelinks]{hyperref}

\title{Soft Soil Creep UMAT --- Test Report}
\author{ConstitutiveModels}
\date{\today}

\begin{document}
\maketitle

\section{Scope}
This report summarises the verification of the Soft Soil Creep (SSC) UMAT
\path{c_models/soft_soil/soft_soil_creep.c}, following \citet{Stolle_1999},
\emph{Time integration of a constitutive law for soft clays}. The tests are in
\path{tests/test_soft_soil_creep.py}; the figures are produced by
\path{docs/ssc_test_report/make_figures.py}, which drives the UMAT along the same
constant rate of strain (CRS) paths as the tests.
Stresses are compression positive; $p' = \tfrac13 \operatorname{tr}\boldsymbol\sigma'$ and
$q = \sqrt{\tfrac32\,\mathbf s:\mathbf s}$.

\section{Model}
The SSC model extends the one-dimensional secondary compression laws of \citet{Buisman_1936},
\citet{Bjerrum_1967} and \citet{Garlanger_1972} to general stress states
\citep{Stolle_1997, Vermeer_1998, Vermeer_1999}. The UMAT implements the formulation of
\citet{Stolle_1999}; equation numbers without further reference refer to that paper.
\begin{itemize}
  \item The creep rate is governed by the equivalent pressure $p_{eq} = p' + q^2/(M^2 p')$
        (Eq.~1), i.e.\ the contours of $p_{eq}$ are the Modified Cam-Clay ellipses of
        \citet{Roscoe_1968}, and creep flow is normal to them (Eq.~7).
  \item Volumetric creep follows the isotache law of Eq.~2, with the pre-consolidation pressure
        $p_p$ growing with the creep strain (Eq.~3).
  \item Elasticity has a bulk modulus $K = p'/A$ and a constant Poisson's ratio (Eq.~6); failure is
        bounded by Mohr--Coulomb with zero dilatancy (Eq.~15).
  \item Stresses are updated implicitly in $p$--$q$ space with a scalar Newton iteration on the creep
        strain increment (Eqs.~11--13), analogous to the scheme of \citet{Borja_1990} for Cam-clay
        plasticity; the consistent tangent is Eq.~9.
\end{itemize}

\begin{table}[h]
\centering
\caption{Material parameters (Table I of the paper).}
\begin{tabular}{cccccccc}
\toprule
$A$ & $B$ & $C$ & $\nu$ & $\tau$ [min] & $M$ & $\varphi$ [$^\circ$] & $M_{mc}$ \\
\midrule
0.016 & 0.090 & 0.004 & 0.25 & 0.70 & 1.29 & 30 & 1.20 \\
\bottomrule
\end{tabular}
\end{table}

\section{Test results}
All 13 unit tests pass (pytest, \today).

\begin{table}[h]
\centering
\small
\begin{tabular}{llc}
\toprule
Test & Checks & Result \\
\midrule
\path{initial_preconsolidation_pressure} & $p_p = \mathrm{OCR}_0\,p_{eq,0}$ on first call & pass \\
\path{elastic_isotropic_response} & $p = p_0\exp(\Delta\varepsilon_v/A)$, elastic tangent & pass \\
\path{zero_time_increment_is_elastic} & no creep for $\Delta t = 0$ & pass \\
\path{tangent_matches_finite_differences_...} & elastic DDSDDE vs.\ finite differences & pass \\
\path{isotropic_creep_..._matches_analytic} & $\varepsilon_c = C\ln(1+t/\tau)$ & pass \\
\path{relaxation_at_constant_strain} & implicit relaxation, decaying rate & pass \\
\path{tangent_is_eq9} & DDSDDE $= \mathbf D^{ec}$ (Eq.~9) & pass \\
\path{oedometer_reaches_k0_line} & CRS oedometer, $K_0 \approx 0.62$ & pass \\
\path{undrained_triaxial_reaches_failure} & CRS undrained, peak $\approx 55$~kPa & pass \\
\path{large_time_steps_match_small_time_steps} & $\Delta t = 21.6$ vs.\ $2.16$ min & pass \\
\path{undrained_strength_increases_...} & rate-dependent peak strength & pass \\
\path{rotation_invariance} & isotropy of the model & pass \\
\path{mohr_coulomb_is_respected_...} & no MC violation on a 3D path & pass \\
\bottomrule
\end{tabular}
\end{table}

\section{CRS oedometer test}
Starting from $p_0 = 10$~kPa with $\mathrm{OCR}_0 = 6$, 20\,\% axial strain is applied in
2160~min (Fig.~2 of the paper). After recompression the path follows the $K_0$ line with
$K_0 = \sigma'_h/\sigma'_v = 0.617$; the final creep strain is $\varepsilon_c = 0.150$ and the
OCR stabilises at 1.22 under the constant strain rate. The 100-step and 1000-step solutions
coincide (Fig.~\ref{fig:oed}).

\begin{figure}[h]
\centering
\begin{subfigure}{0.49\textwidth}\includegraphics[width=\textwidth]{figures/oedometer_eps_q.pdf}\caption{$\varepsilon_1$--$q$ (log scale)}\end{subfigure}
\begin{subfigure}{0.49\textwidth}\includegraphics[width=\textwidth]{figures/oedometer_p_q.pdf}\caption{$p'$--$q$}\end{subfigure}
\caption{CRS oedometer test, $p_0 = 10$~kPa, $\mathrm{OCR}_0 = 6$.}
\label{fig:oed}
\end{figure}

\section{CRS undrained triaxial test}
A normally consolidated sample at $p_0 = 100$~kPa is sheared isochorically to 20\,\% axial strain
in 2160~min (Fig.~3 of the paper). Creep generates excess pore pressure, so $p'$ decreases
monotonically. $q$ peaks at 55.7~kPa (paper: about 55~kPa), after which the path reaches the
Mohr--Coulomb envelope $q = M_{mc}\,p'$ and softens along it to $q = 46.7$~kPa
(Fig.~\ref{fig:und}). Large and small time steps give the same response.

\begin{figure}[h]
\centering
\begin{subfigure}{0.49\textwidth}\includegraphics[width=\textwidth]{figures/undrained_eps_q.pdf}\caption{$\varepsilon_1$--$q$ (log scale)}\end{subfigure}
\begin{subfigure}{0.49\textwidth}\includegraphics[width=\textwidth]{figures/undrained_p_q.pdf}\caption{$p'$--$q$}\end{subfigure}
\caption{CRS undrained triaxial test, normally consolidated, $p_0 = 100$~kPa.}
\label{fig:und}
\end{figure}

\section{Strain-rate dependence}
The undrained test is repeated with test durations of 216, 2160 and 21600~min (100 steps each).
The peak strength increases with strain rate (Table~\ref{tab:rate}, Fig.~\ref{fig:rate}),
by roughly 9\,\% per decade of strain rate.

\begin{table}[h]
\centering
\caption{Peak and final deviator stress of the undrained CRS tests.}
\label{tab:rate}
\begin{tabular}{rcccc}
\toprule
$t_{end}$ [min] & $q_{peak}$ [kPa] & $p'$ at peak [kPa] & $q_{end}$ [kPa] & $p'_{end}$ [kPa] \\
\midrule
21600 & 51.06 & 45.45 & 42.82 & 35.68 \\
2160  & 55.70 & 49.49 & 46.71 & 38.92 \\
216   & 60.76 & 53.86 & 50.94 & 42.45 \\
\bottomrule
\end{tabular}
\end{table}

\begin{figure}[h]
\centering
\begin{subfigure}{0.49\textwidth}\includegraphics[width=\textwidth]{figures/rate_eps_q.pdf}\caption{$\varepsilon_1$--$q$ (log scale)}\end{subfigure}
\begin{subfigure}{0.49\textwidth}\includegraphics[width=\textwidth]{figures/rate_p_q.pdf}\caption{$p'$--$q$}\end{subfigure}
\caption{Undrained CRS tests at three strain rates.}
\label{fig:rate}
\end{figure}

\bibliographystyle{plainnat}
\bibliography{../refs}

\end{document}
